\section{The \mDS{} Run Lost 29\% of Its Rows}
\label{app:truncation}

Roughly 1{,}000 rows into the \mDS{} run, the pinned provider endpoint
began returning HTTP~429. The harness had no retry for that status, so each
throttled request was written to the results file as a terminal error: 466
of 468 error rows are rate limits. Error rate was 0.2\% over the first
1{,}000 rows, 31\% over the next 200, and 100\% thereafter. The sweep then
finished early, because a throttled request fails instantly, and exited
cleanly.

The loss costs power, not validity, and the evidence for that is specific:

\begin{itemize}
\itemsep2pt
\item Errors are near-uniform across arms: 113, 116, 118 and 121 of 401
      rows. No arm is preferentially damaged.
\item 114 of the 122 lost tasks lost \emph{all four} arms together, because
      the runner schedules a task's arms close in time. Independent per-arm
      failure would have left roughly 100 complete tasks; 279 survived.
      Missingness is a property of \emph{when} a task ran.
\item Figures on the 279 tasks scoreable in all four arms (the
      complete-case set, used for \mDS{} throughout) and per-arm figures
      over each arm's own denominator agree to within one point.
\end{itemize}

We therefore report \mDS{} on the 279 tasks scoreable in all four arms,
with intervals roughly 20\% wider than the other runs'. One residual caveat:
the lost tasks are slightly enriched for hard tasks (87\% hard against 83\%
overall), so \mDS{}'s absolute accuracies are marginally optimistic. Every
arm is affected identically. The 122 tasks are re-runnable, but only on a
different provider than the other 279, which would be its own confound;
leaving them out is cleaner.

One lesson generalises beyond this paper. Pinning a single provider endpoint
with fallbacks disabled is correct for validity, since it holds
quantization, price and throughput fixed across arms, but it transfers
responsibility for transient failures to the harness, and here nothing
looked wrong until the results were read. Benchmark harnesses that pin
endpoints need a rate-of-progress check, not just a liveness check. The
harness now retries 429.
